\newpage
\begin{center}
  {\fontsize{18pt}{21pt}\bodyfont\bfseries\color{darktext} Full Curriculum Vitae}
\end{center}
\vspace{-1mm}

\cvsection{Professional Profile}
\begin{cvparagraph}
\textbf{Research profile:} Geoscientist specializing in structural geology, geomechanics, and multiscale numerical modelling, with more than 15 years of experience in research, university teaching, fieldwork, and hands-on student training. His research addresses geodynamic, geothermal, and geomechanical processes in Andean systems and energy reservoirs, integrating two- and three-dimensional models to study subduction zones, crustal deformation, volcanic systems with geothermal potential, and natural and induced fractures, with emphasis on the Vaca Muerta Formation.

\textbf{Teaching, outreach, and knowledge transfer:} Teaching experience in Geological Surveying, Geostatistics, and Introduction to Geology, including practical work in geological mapping, spatial analysis, geological interpretation, geospatial tools, UAVs, and quantitative methods. Participation in public science communication, STEAM education, and community initiatives addressing territorial and urban-environmental issues and access to public services.
\end{cvparagraph}

\cvsection{Academic and Postdoctoral Training}
\begin{cventries}
\cventry{CONICET and UPC}{Postdoctoral Fellowship}{Argentina/Spain}{2015--2017}{\begin{cvitems}\item{Research topic: thermal evolution of oceanic lithosphere. Supervisors: Dr. Ernesto Cristallini (CONICET--UBA) and Dr. Sergio Zlotnik (UPC).}\end{cvitems}}
\cventry{Stanford University}{Visiting Researcher}{USA}{2013}{\begin{cvitems}\item{Courses in Structural Geology and MATLAB applied to Geosciences.}\end{cvitems}}
\cventry{University of Buenos Aires}{PhD in Geological Sciences}{Argentina}{2010--2015}{\begin{cvitems}\item{Thesis: \emph{Prediction of fracturing on irregular geological surfaces using static methods}. Supervisor: Dr. Ernesto Cristallini.}\end{cvitems}}
\cventry{University of Buenos Aires}{BSc in Geological Sciences}{Argentina}{2005--2010}{\begin{cvitems}\item{Thesis: \emph{Geology and structure of the Tres Cruces compound anticline, Jujuy, Argentina}. Supervisor: Dr. Ernesto Cristallini.}\end{cvitems}}
\end{cventries}

\cvsection{Appointments and Professional Experience}
\cvsubsection{Current appointments}
\begin{cventries}
\cventry{CONICET--IDEAN}{Adjunct Researcher (Investigador Adjunto)}{Argentina}{2018--present}{\begin{cvitems}\item{Research areas: fracture prediction, geodynamic modelling, and geothermal-potential modelling. CONICET research career from 2018; promoted to Investigador Adjunto, November 2023.}\end{cvitems}}
\cventry{Department of Geological Sciences, University of Buenos Aires}{Associate Professor}{Argentina}{May 2026--present}{\begin{cvitems}\item{Courses: Geostatistics and Geoinformatics.}\end{cvitems}}
\cventry{IDEAN, CONICET--UBA}{Institutional Coordinator}{Argentina}{Sep 2025--present}{\begin{cvitems}\item{Council Member since April 2024.}\end{cvitems}}
\end{cventries}

\cvsubsection{Previous teaching appointment}
\begin{cventries}
\cventry{Department of Geological Sciences, University of Buenos Aires}{Teaching Assistant}{Argentina}{2013--2026}{\begin{cvitems}\item{Courses: Introduction to Geology, Geological Surveying, and Geostatistics.}\end{cvitems}}
\end{cventries}

\cvsubsection{Professional courses and field activities}
\begin{cventries}
\cventry{Instructor in charge}{Geology of Unconventional Reservoirs}{Argentina}{2025}{\begin{cvitems}\item{Professional training course focused on the geological characterization of shale reservoirs.}\end{cvitems}}
\cventry{Field instructor}{Advanced Research Field Trip}{Argentina}{2025}{\begin{cvitems}\item{Organizer and instructor for \emph{Perspectives on Vaca Muerta: subsurface attributes analysed from surface exposures}.}\end{cvitems}}
\end{cventries}

\newpage
\cvsubsection{Research visits}
\begin{cventries}
\cventry{GEO3BCN--CSIC}{Visiting Researcher}{Spain}{May--Sep 2025}{\begin{cvitems}\item{Thermomechanical modelling of volcanic zones for geothermal-potential assessment. Host: Dr Ivone Jiménez-Munt.}\end{cvitems}}
\cventry{GEO3BCN--CSIC}{Visiting Researcher}{Spain}{Aug--Sep 2024}{\begin{cvitems}\item{Development of coupled numerical models for volcanic and geothermal systems.}\end{cvitems}}
\cventry{BarcelonaTech (UPC)}{Visiting Researcher}{Spain}{2015, 2017, 2019, 2022}{\begin{cvitems}\item{Research stays at the Numerical Computation Laboratory.}\end{cvitems}}
\end{cventries}

\cvsection{Research Projects}
\begingroup
\renewcommand{\arraystretch}{1.12}
\renewcommand*{\honorpositionstyle}[1]{{\fontsize{9pt}{10.8pt}\bodyfont\mdseries\color{darktext} #1}}
\begin{cvhonors}
\cvhonor{Geothermal potential of volcanic zones (Principal Investigator).}{AECT-2025-1-0009}{Spain}{2025}
\cvhonor{Structural controls in volcanic systems: magmatic intrusions and damage-zone development associated with hydrocarbon generation (Co-PI).}{Williams Foundation}{Argentina}{2025}
\cvhonor{Assessment of geothermal potential in volcanic zones using thermomechanical numerical models (PI).}{LINCGL2024}{Spain}{2024}
\cvhonor{Geodynamic modelling of subduction zones and mantle plumes (PI).}{AECT-2024-2-0027}{Spain}{2024}
\cvhonor{Geodynamics of the Adria microplate, from mantle to surface (Co-PI).}{GEO3BCN}{Spain}{2023}
\cvhonor{Geodynamic modelling of subduction zones (PI).}{AECT-2022-3-0023}{Spain}{2022}
\cvhonor{Landscape dynamics and geological hazards associated with dams in mountain regions (Co-PI).}{PICT-2019-04011}{Argentina}{2020}
\cvhonor{Analogue and numerical modelling: contributions to tectonic and structural problems (Co-PI).}{PICT-2019-00997}{Argentina}{2020}
\cvhonor{Natural fracturing and first-order structures: predictive models applied to Cerro Domuyo (Co-PI).}{UBACyT}{Argentina}{2018}
\cvhonor{Geological evolution of the Andes and its economic and environmental impact (Co-PI).}{PUE 22920160100051}{Argentina}{2017}
\cvhonor{Numerical modelling of the thermal evolution of oceanic lithosphere (PI).}{PICT 2015-1226}{Argentina}{2016}
\cvhonor{Application of analogue and numerical models to tectonic-structural problems (Co-PI).}{PICT-2016-1407}{Argentina}{2016}
\end{cvhonors}
\endgroup

\cvsection{Research Supervision}
\cvsubsection{Researcher supervision}
\begin{cventries}
\cventry{Assistant Researcher}{Plotek, Berenice}{Argentina}{Since Aug. 2025}{\begin{cvitems}\item{Supervisor: Likerman, J.; co-supervisor: Cristallini, E.}\end{cvitems}}
\end{cventries}

\cvsubsection{Postdoctoral supervision}
\begin{cventries}
\cventry{Fault-propagation folds: a mechanical approach using finite-element numerical modelling}{Plotek, Berenice}{Argentina}{2024}{\begin{cvitems}\item{Supervisors: Cristallini, E. and Likerman, J.}\end{cvitems}}
\end{cventries}

\cvsubsection{Doctoral supervision}
\begin{cventries}
\cventry{Natural fracturing in the Vaca Muerta Formation: characterization and modelling}{Correa Luna, Clara}{Argentina}{2023}{\begin{cvitems}\item{Supervisors: Yagupsky, D. and Likerman, J.}\end{cvitems}}
\end{cventries}

\newpage
\cvsection{Undergraduate Thesis Supervision}
\begin{cventries}
\cventry{Tectonic inversion of a metamorphic complex in the Deseado Massif: insights from analogue modelling}{Chandía Sepúlveda, Alexander Ignacio}{Argentina}{In progress}{\begin{cvitems}\item{Supervisors: Likerman, J. and Navarrete, C.}\end{cvitems}}
\cventry{Quaternary deformation in the Fiambalá Valley, Catamarca}{De Rosa, Tomás}{Argentina}{In progress}{\begin{cvitems}\item{Supervisors: Likerman, J. and Sagripanti, L.}\end{cvitems}}
\cventry{Geothermal-electric potential of the Cerro Blanco Volcanic Complex: influence of rheology and structure through numerical models}{Blanco, Lucía}{Argentina}{In progress}{\begin{cvitems}\item{Supervisors: Likerman, J. and Lamberti, C.}\end{cvitems}}
\cventry{Structural evolution of an inverted metamorphic complex in the Deseado Massif: insights from numerical modelling}{Celis Santisteban, Celso Angel}{Argentina}{In progress}{\begin{cvitems}\item{Supervisors: Likerman, J. and Navarrete, C.}\end{cvitems}}
\cventry{Assessment of the geothermal-electric potential of the Cerro Blanco Volcanic Complex using numerical modelling}{Ruiz, Luciana}{Argentina}{2026}{\begin{cvitems}\item{Supervisors: Lamberti, C. and Likerman, J.}\end{cvitems}}
\cventry{Quaternary deformation in the Chaschuil Valley, western Sierra de Las Planchadas, Catamarca}{Gramajo, Pilar}{Argentina}{2025}{\begin{cvitems}\item{Supervisors: Likerman, J. and Sagripanti, L.}\end{cvitems}}
\cventry{Geomorphology of the Río Grande valley between Bardas Blancas and Portezuelo del Viento, Mendoza}{Acosta, Luana}{Argentina}{2023}{\begin{cvitems}\item{Supervisors: Winocur, D. and Likerman, J.}\end{cvitems}}
\cventry{Geology of the eastern slope of Cerro Perito Moreno, El Bolsón, Río Negro}{Mazzini, Martín}{Argentina}{2022}{\begin{cvitems}\item{Supervisors: Likerman, J. and Tobal, J.}\end{cvitems}}
\cventry{Structure and neotectonics of the Tabladitas area, Jujuy, Argentina}{Relañez, Gastón}{Argentina}{2019}{\begin{cvitems}\item{Supervisors: Yagupsky, D. and Likerman, J.}\end{cvitems}}
\cventry{Geology, structure, and fracturing of the Lajas Formation in the Covunco anticline, Neuquén}{Ventisky, Esteban}{Argentina}{2017}{\begin{cvitems}\item{Supervisors: Likerman, J. and Tobal, J.}\end{cvitems}}
\cventry{Structural geology and fracturing in the Tres Cruces Basin, Jujuy}{Correa Luna, Clara}{Argentina}{2015}{\begin{cvitems}\item{Supervisors: Yagupsky, D. and Likerman, J.}\end{cvitems}}
\end{cventries}

\newpage
\printbibsection{article}{Peer-Reviewed Publications}

\cvsection{Fellowships and Distinctions}
\begin{cvhonors}
\cvhonor{Research Fellowship in Japan}{Kobe University}{JSPS}{2027}
\cvhonor{SZNet 2025 Chile Field Trip}{National Science Foundation}{SZ4D}{2025}
\cvhonor{Academic Mobility Fellowship among Partner Institutions}{Ibero-American University Association for Postgraduate Studies}{AUIP}{2022}
\cvhonor{International Research Stay}{National Scientific and Technical Research Council}{CONICET}{2016}
\cvhonor{Postdoctoral Fellowship}{National Scientific and Technical Research Council}{CONICET}{2015}
\cvhonor{Doctoral Research Stay at Stanford University}{Fulbright -- Bunge \& Born Foundation}{Doctoral fellowship}{2013}
\cvhonor{Doctoral Fellowship Type II}{National Scientific and Technical Research Council}{CONICET}{2013}
\cvhonor{Doctoral Fellowship Type I}{National Scientific and Technical Research Council}{CONICET}{2010}
\end{cvhonors}

\cvsection{Skills}
\descriptionstyle{\textbf{Programming:} \LaTeX{}, Python, R, MATLAB, HTML5.}\par
\vspace{0.8mm}
\descriptionstyle{\textbf{Modelling tools:} Underworld2, LaMEM, GOLEM, ASPECT; analogue modelling, experimental design, material properties, PIV, and post-processing; HPC workflows.}\par
\vspace{0.8mm}
\descriptionstyle{\textbf{Spatial analysis:} QGIS, GMT, PostGIS; seismic interpretation, geological mapping, and UAV-assisted field methods.}\par
\vspace{0.8mm}
\descriptionstyle{\textbf{Fieldwork:} more than 15 years of experience, extensive knowledge of the Andes, campaign planning and logistics, and team leadership.}\par
\vspace{0.8mm}
\descriptionstyle{\textbf{Languages:} Spanish (native), English (fluent, C2), Italian (A2), Catalan (A2).}
